\section{Principe de Mach, propagation bidirectionnelle et réponse gravitationnelle collective}
\label{chap:mach-onda-bidireccional-rev10}
\index{Mach, principe de}
\index{onde!avancée et retardée}
\index{Doppler!compression et expansion}
\index{gravitation!réponse collective}

\subsection{Dépendance globale de la réponse locale}

Soit \(\mathcal X_\infty\) l’espace des sections compatibles, soit \(\operatorname{res}_v:\mathcal X_\infty\to X_v\) la restriction à un voisinage local et soit \(b:\mathcal X_\infty\to B_\infty\) la trace globale de frontière. Une réponse inertielle locale \(I_v:\mathcal X_\infty\to V_v\) est \emph{machienne} lorsqu’il existe \(\overline I_v:\operatorname{im}(b)\to V_v\) tel que
\begin{equation}
 I_v=\overline I_v\circ b,
 \label{eq:factorizacion-mach-rev10}
\end{equation}
mais \(I_v\) ne se factorise pas seulement par \(\operatorname{res}_v\). En particulier, il peut exister \(x,x'\) avec
\[
 \operatorname{res}_v(x)=\operatorname{res}_v(x'),\qquad
 b(x)\ne b(x'),\qquad I_v(x)\ne I_v(x').
\]
La définition formalise le principe de Mach : la réponse locale dépend de la classe globale de fermeture et non d’un espace absolu ajouté \cite{MachMechanics}. C’est le même schéma local--global qui apparaît dans les cylindres de nombre, dans l’holonomie et dans le transport de mémoire.

\begin{theorem}[Critère de factorisation de la réponse par la totalité]
\label{thm:criterio-factorizacion-machiana-rev2}
Il existe une unique application \(\overline I_v:\operatorname{im}(b)\to V_v\) satisfaisant \(I_v=\overline I_v\circ b\) si et seulement si \(I_v\) est constante sur chaque fibre de \(b\) :
\[
 b(x)=b(y)\quad\Longrightarrow\quad I_v(x)=I_v(y).
\]
Si \(b\) est surjective, la réponse induite est définie de manière unique sur tout \(B_\infty\).
\end{theorem}

\begin{proof}
La nécessité résulte de l’application de \(\overline I_v\) à l’égalité \(b(x)=b(y)\). Pour la suffisance, étant donné \(\beta\in\operatorname{im}(b)\), définissons \(\overline I_v(\beta)=I_v(x)\) pour tout \(x\in b^{-1}(\beta)\). La constance sur les fibres rend le choix du représentant sans effet ; la factorisation fixe ensuite la valeur de \(\overline I_v\) en chaque point de son domaine et garantit l’unicité.
\end{proof}

La relation entre état et coordonnée est fixée par
\[
 [x]_{\mathfrak F}\ne x,
 \qquad
 q(x)\ne x,
 \qquad
 x\in\mathcal X_\infty.
\]
L’état conserve position, feuillet, orientation, mode additif ou multiplicatif, quotient, retenue, frontière, parcours et mémoire ; l’observable n’exprime qu’une coordonnée. Le repère local correspondant n’est pas ajouté de l’extérieur, mais induit par restriction de la totalité compatible et de sa connexion :
\begin{equation}
 \mathfrak F_v(x)
 :=\operatorname{Res}_v\bigl(x,
          \nabla_{\rm TPK}\bigr).
 \label{eq:marco-local-totalidad-rev10}
\end{equation}
Ici, \(\mathcal X_\infty\) ne désigne pas une somme verbale de tout l’univers, mais le système des sections compatibles avec orientation, frontière et mémoire. L’autoréférence est coinductive, non circulaire : chaque section est reconnue par la compatibilité de toutes ses troncatures.

Cette formulation organise les régimes inertiels et accélérés au moyen d’une unique connexion endogène. Pour une courbe \(\gamma\), l’accélération covariante est
\[
 a_\gamma:=\nabla_{\dot\gamma}\dot\gamma,
 \qquad
 a_\gamma=0
 \quad\Longleftrightarrow\quad
 \gamma\ \text{est une géodésique affine de }\nabla.
\]
Si \((e_a)\) est une tétrade le long de \(\gamma\), la condition
\[
 \nabla_{\dot\gamma}e_a=0
\]
définit un repère parallèle et non tournant relativement à la connexion. La rotation d’un repère général se mesure par ses coefficients de connexion le long de \(\gamma\), ou par son écart au transport de Fermi--Walker lorsque celui-ci est la référence cinématique choisie. Pour un champ arbitraire \(v\), \(D_\gamma v=0\) signifie uniquement transport parallèle ; sa négation n’équivaut pas à elle seule à une accélération ou à une rotation.

La différence physique est ainsi mesurée au moyen de l’accélération covariante, de la rotation du repère, de la torsion, de la courbure et de l’holonomie, sans introduire de repère privilégié extérieur. En ce sens, chaque section persistante est \emph{auto-inertielle} : elle porte le repère induit par sa propre insertion dans la totalité HMT.

\begin{proposition}[Auto-inertie, retour et holonomie]
\label{prop:autoinercia-holonomia-rev2}
Soit \(E\to M\) le fibré de réalisation, soit \(\gamma\) un lacet basé en \(x_0\) et soit \(\operatorname{Hol}_{\nabla}(\gamma):E_{x_0}\to E_{x_0}\) l’holonomie de la connexion induite. Une donnée \(s_0\in E_{x_0}\) transportée parallèlement revient exactement si et seulement si
\[
 s_0\in\operatorname{Fix}
 \bigl(\operatorname{Hol}_{\nabla}(\gamma)\bigr).
\]
Si l’état généalogique complet est réalisé dans \(E_{x_0}\) par une application injective sur le secteur comparé, cette condition équivaut aussi au retour de l’état complet. Sans cette injectivité, la conclusion porte exclusivement sur la donnée dans la fibre. Si la coordonnée visible revient et que \(s_0\) reste hors du sous-espace fixe, il y a retour de phase avec changement de l’état transporté.
\end{proposition}

\begin{proof}
Le transport autour du lacet agit sur \(E_{x_0}\) au moyen de \(\operatorname{Hol}_{\nabla}(\gamma)\). L’égalité entre les données finale et initiale équivaut, par définition, au fait que \(s_0\) est un vecteur fixe de l’opérateur. La dernière équivalence s’obtient en appliquant l’injectivité de la réalisation au secteur déclaré.
\end{proof}

\subsection{\(2\) orientations d’une même propagation}

Pour un parcours \(\gamma\) entre émission et absorption, l’inversion \(\gamma^\vee\) échange les extrémités et conserve la longueur de la trajectoire. Dans la réalisation photonique de \cref{chap:realizaciones-nula-material}, le couple \((\gamma,\gamma^\vee)\) représente les orientations retardée et avancée d’un même transport nul. Un ruban de Möbius conserve cette mémoire d’orientation : un tour échange les feuillets et deux tours rétablissent la section orientée.

Cette structure permet de formuler rigoureusement l’intuition d’une onde pilote bidirectionnelle. L’état matériel n’est pas identifié à une onde supplémentaire qui le pousse de l’extérieur ; la section et sa phase sont deux lectures du même transport enrichi. L’identification concrète aux propagateurs de Green avancé et retardé exige néanmoins un entrelaceur préservant l’action, la causalité et le produit scalaire.

\begin{proposition}[Conjugaison d'orientation et réversion du propagateur]
\label{prop:conjugacion-propagador-bidireccional-rev2}
Soit \(H=H^*\) un opérateur autoadjoint, éventuellement non borné, qui satisfait
\[
 HJ_E=J_EH.
\]
Soit \(C\) une involution orthogonale de la réalification quantique qui préserve \(D(H)\) et satisfait, sur ce domaine,
\[
 CJ_EC^{-1}=-J_E,
 \qquad CHC^{-1}=H.
\]
Alors
\[
 U(t):=\exp(-J_EHt/\hbar_{\rm int})
\]
est un groupe unitaire à un paramètre et l’on a
\[
 \boxed{CU(t)C^{-1}=U(-t)}.
\]
\end{proposition}

\begin{proof}
L’opérateur \(A=-J_EH/\hbar_{\rm int}\) est anti-autoadjoint et engendre \(U(t)\) par le théorème de Stone. Les hypothèses de domaine et de conjugaison donnent \(CAC^{-1}=-A\) ; le calcul fonctionnel, ou de manière équivalente l’unicité du groupe engendré, implique \(C\exp(tA)C^{-1}=\exp(-tA)\). Si \(H\) est borné, la même identité s’obtient terme à terme à partir de la série exponentielle.
\end{proof}

L’involution de mot \(C_M=-\operatorname{rev}\) fournit la donnée généalogique d’orientation. Lorsque sa réalisation satisfait les hypothèses de la proposition, les branches particule et antiparticule sont représentées par des propagateurs de sens temporels conjugués : la géométrie du ruban fixe l’orientation et le Hamiltonien fixe la dynamique.

La lecture cinématique de compression et d’expansion est fixée, pour \(|\beta|<1\), par la rapidité \(\zeta=\operatorname{artanh}\beta\). Pour une fréquence propre \(\omega_0\), les deux orientations Doppler sont
\begin{equation}
 \omega_{\rm comp}=\omega_0e^{\zeta}
 =\omega_0\sqrt{\frac{1+\beta}{1-\beta}},\qquad
 \omega_{\rm exp}=\omega_0e^{-\zeta}
 =\omega_0\sqrt{\frac{1-\beta}{1+\beta}},
 \label{eq:doppler-bidireccional-rev10}
\end{equation}
et satisfont
\[
 \omega_{\rm comp}\omega_{\rm exp}=\omega_0^2.
\]
La réciprocité de \eqref{eq:doppler-bidireccional-rev10} est une réalisation continue du schéma bicouche : une orientation comprime et sa conjuguée dilate, tandis que leur produit conserve l’échelle centrale. Cela n’identifie pas le paramètre Doppler à \(C^*\), à \(R_{\rm act}\) ni à un coefficient de masse ; les applications entre ces coordonnées doivent être construites séparément.

\subsection{Orbite elliptique, action et \(3\) invariants J}

L’orbite de Kepler, son relèvement de Levi--Civita et la condition de Sommerfeld organisent la phase sur un cycle elliptique. Dans cette branche,
\[
 J_\gamma=\oint_\gamma p\,\mathrm dq
\]
est l’action orbitale. Ce symbole ne se confond ni avec le caractère modulaire \(J_{\rm mod}=j-744\), qui relève de Moonshine, ni avec l’invariant de Jarlskog \(J_{\rm CKM}\), qui mesure la violation CP. Les trois objets partagent une propriété abstraite --- ce sont des invariants obtenus après fermeture d’une composition --- mais possèdent des domaines, des unités et des théorèmes distincts.

La condition
\[
 \varepsilon_\gamma
 \exp\!\left(\frac{iJ_\gamma}{\hbar_{\rm int}}\right)=1
\]
explique pourquoi l’holonomie de signe décale d’une demi-unité la règle de Sommerfeld. Avec \(\operatorname{Area}(\theta)=\hbar_{\rm HMT}^{\rm ret}\theta\), elle rend visibles la fermeture \(2\pi\) et le double tour \(4\pi\) : la première achève le cycle scalaire ; le second rétablit une section de spin ou de Möbius qui a changé de feuillet après un seul tour.

\subsection{Courbure propagative comme mode collectif de la géométrie}

La variable gravitationnelle primaire est le transport : cotétrade, connexion, courbure, torsion et holonomie. Une onde gravitationnelle est une variation propagative de la courbure ou de sa classe d’holonomie. Le type d’état ne nécessite pas de coordonnée élémentaire supplémentaire de spin deux. Un quantum d’hélicités \(\pm2\), s’il apparaît après réduction des contraintes et quantification des fluctuations, est une excitation collective du secteur géométrique et non un constituant primitif.

La lecture machienne renforce cette organisation : la perturbation locale mesure un changement de la section globale. La composante qui ne demeure pas dans la courbure observable ne peut être représentée comme mémoire locale de transport ou contorsion que lorsque le diagramme entre TPK et Cartan commute ; la mémoire discrète n’est pas identifiée à la torsion par une égalité nominale. Sous cette condition, l’onde peut se propager sans postuler un graviton fondamental et la réponse locale conserve l’information de la fermeture globale.

\subsection{Hypothèses métriques et conditions aux limites de la réalisation machienne}

La sémantique temporelle du TRIT attribue au retour la mémoire du passé, au seuil l’actualisation présente et à l’ouverture le prolongement futur. La réalisation elliptique, euclidienne et hyperbolique traduit ces trois régimes en fermeture, frontière et divergence. Cela fournit une lecture géométrique de la matière comme persistance enroulée et de la propagation comme ouverture hyperbolique.

Interpréter une particule fermée comme une section bidimensionnelle effondrée, un électron comme la fermeture de deux modes nuls, ou un trou noir comme une inversion globale de feuillets constitue des conséquences théoriques compatibles avec cette architecture. Ce ne sont pas des identités purement combinatoires : pour acquérir un statut physique, elles doivent réaliser une métrique, résoudre les équations de champ et reproduire l’horizon, la causalité et les conditions de stabilité correspondantes. Cette séparation préserve la puissance explicative du modèle sans transformer une interprétation géométrique en preuve déjà obtenue.

\begin{statusbox}
La factorisation machienne, la monodromie de Möbius, la réciprocité Doppler et la condition de phase sont des constructions mathématiques explicites dans leurs domaines déclarés. Leur identification conjointe à des ondes pilotes physiques, à des trous noirs ou à la cosmologie constitue une réalisation théorique qui doit satisfaire les entrelaceurs et les équations indiqués.
\end{statusbox}
